\chapter{Future Work}

There is a plethora of features and capabilities that can be
developed on this project. Following the CAMARAaaS philosophy,
the adoption of the 3GPP Common API Framework (CAPIF) is the next logical
step. This allows for the controlled exposure and consumption of
the APIs \cite{CAMARAaaS}. CAPIF defines a set of core components
\cite{CAPIF}:

\begin{itemize} \item \textbf{API Invoker}: A software module
	      that is usually provided by 3rd party but is present in
	      trusted operator Public Land

	\item \textbf{CAPIF Core Function}: The module responsible
	      for authenticating the API Invoker before the API services
	      are reached. It also publishes and stores necessary API
	      information, and interacts with other needed CAPIF Core
	      Functions.

	\item \textbf{API Exposing Function}: The entity which
	      provides the service communication entry point for the
	      service APIs.

	\item \textbf{API Publishing Function}: The entity that
	      enables the API provider to publish the Service APIs
	      information in order to enable the discovery of APIs by the
	      API invoker.

	\item \textbf{API Management Function}: The entity which
	      enables the API provider to perform administration of the
	      service APIs.

\end{itemize}

It also defines a plethora of requirements that must be respected, namely in
security, charging, logging, and policy configuration concerns.


It is also of extreme importance to develop the remaining CAMARA APIs, in order
to achieve full functionality of the gateway. This project consisted on the
implementation of a subset of the CAMARA APIs, not the full suite, in part due
to time constraints and network limitations. This would imply the development
of an additional 26 APIs.

Due to the research nature of IT's 5G deployment, it is not unreasonable to
expect a lot of auxiliary work once again, in order to implement said APIs, so
that is something that would also need to be performed.

Lastly, due to the way our architecture is laid out, it would be necessary to
develop more drivers to interact with different operators and 5G Cores.
Currently, our only interaction is with Huawei's 5G Core present in the IT
infrastructure, however, development effort would need to be performed to allow
our gateway to interact, for example, with CumuCore's offerings.
